\chapter*{Was Nullprodukte verraten}
\hypertarget{nullproducts.reading}{}
\addcontentsline{toc}{chapter}{Was Nullprodukte verraten}

Eine Halbgruppe besteht aus Elementen, die man assoziativ multiplizieren
kann. Man kann aber auch ihre nichtleeren Teilmengen multiplizieren:
Man nimmt aus jeder der beiden Mengen ein Element, bildet dessen Produkt
und sammelt alle so entstehenden Werte. Für eine Halbgruppe \(S\) ist also
\[
 AB=\{ab\mid a\in A,\ b\in B\}\qquad(A,B\in\PowNE S).
\]
Beispielsweise enthält \(\{a,b\}\{u,v\}\) die Werte \(au,av,bu,bv\).
Einige dieser Werte können gleich sein. Das Ergebnis muss daher keine
vier Elemente enthalten. Die Menge aller nichtleeren Teilmengen mit
dieser Multiplikation heißt die \emph{Potenzhalbgruppe} von \(S\).

Wie viel von der ursprünglichen Halbgruppe lässt sich aus dieser
größeren Multiplikationstafel erkennen? Hier untersuchen wir eine
endliche Halbgruppe \(S\) mit Nullelement \(0\), in der genau zwei
verschiedene Produktwerte vorkommen:
\[
 S^2=\{xy\mid x,y\in S\}=\{0,c\},\qquad c\ne0.
\]
\(S^2\) bezeichnet hier die Menge der Produktwerte, nicht das kartesische
Produkt. Jedes Produkt ist also entweder \(0\) oder \(c\), und beide
Werte treten auf. Das Nullelement erfüllt \(0x=x0=0\) für jedes \(x\in S\).

\medskip\noindent\textbf{Was genau soll bewiesen werden?}
Sei \(T\) eine weitere Halbgruppe. Gegeben sei ein Isomorphismus
zwischen den beiden Potenzhalbgruppen, also eine bijektive Abbildung
\[
 \Phi\colon\PowNE S\longrightarrow\PowNE T,\qquad
 \Phi(AB)=\Phi(A)\Phi(B).
\]
\(T\) bezeichnet im ganzen Text diese zweite Halbgruppe. Ihre
Multiplikation schreiben wir ebenfalls durch Nebeneinandersetzen.
Gesucht ist eine Bijektion \(f\colon S\to T\) mit
\[
 f(xy)=f(x)f(y)\qquad(x,y\in S).
\]
Diese Eigenschaft macht \(f\) zu einem Halbgruppenisomorphismus:
Die beiden Halbgruppen haben dieselbe Multiplikation bis auf Umbenennung
ihrer Elemente. Dies ist der Satz
\FormulaRefAuto{FiniteTwoProductZeroReconstruction}[thm] aus Band~37.
Die Endlichkeit und ein Nullelement von \(T\) werden nicht zusätzlich
vorausgesetzt; wir werden beides herleiten.

Die beiden Abbildungen haben unterschiedliche Aufgaben. \(\Phi\) ist
bereits gegeben und ordnet einer \emph{nichtleeren Teilmenge} von \(S\)
eine nichtleere Teilmenge von \(T\) zu. Dagegen ist \(f\) noch zu
konstruieren und soll \emph{einzelne Grundelemente} zuordnen.
Aus dem gegebenen \(\Phi\) wird erst im Laufe des Beweises ein passendes
\(f\) gewonnen.

\section*{Warum wir die Einermengen nicht einfach ablesen können}

Für Einermengen gilt stets \(\{x\}\{y\}=\{xy\}\). Würde \(\Phi\)
die Einermengen bijektiv auf die Einermengen abbilden, könnten wir
\(\Phi(\{x\})=\{f(x)\}\) setzen und die gesuchte Zuordnung unmittelbar
ablesen. Aber ein Isomorphismus der Potenzhalbgruppen muss nur die
Multiplikation erhalten. Dass er auch Einermengen, Inklusionen oder
Vereinigungen erhält, gehört nicht zu den Voraussetzungen.

Ein kleines Beispiel zeigt das Problem. Auf \(N=\{0,a\}\) seien alle
Produkte gleich \(0\). Die Potenzhalbgruppe hat die drei Elemente
\(\{0\},\{a\},\{0,a\}\), und jedes Produkt zweier dieser Elemente ist
\(\{0\}\). Vertauschen wir \(\{a\}\) und \(\{0,a\}\) und lassen
\(\{0\}\) fest, bleibt die gesamte Multiplikation erhalten:
Das Bild jedes Produkts und das Produkt der beiden Bilder sind
weiterhin \(\{0\}\). Dieser Automorphismus vertauscht eine Einermenge
mit einer Zweiermenge.
Dieses Beispiel hat nur einen Produktwert. Es zeigt, warum die
allgemeinen Isomorphismusbedingungen allein das Erkennen von
Einermengen nicht rechtfertigen.

Wir werden deshalb nicht behaupten, dass \(\Phi\) selbst die gesuchte
Zuordnung der Grundelemente enthält. Stattdessen vergleichen wir,
welche Mengen miteinander das Produkt \(\{0\}\) ergeben. Mengen mit
demselben Verhalten werden zusammengefasst. Diese Familien heißen
später \emph{Nullproduktprofile}. Aus ihrer Ordnung und aus endlichen
Anzahlen bestimmen wir, wie viele Grundelemente zu jedem Profil gehören.
Dann können wir die Grundelemente passend zuordnen.

\medskip\noindent\textbf{Was ist ein Nulltest?}
Wir halten eine nichtleere Menge \(A\subseteq S\) fest und wählen
eine nichtleere \emph{Testmenge} \(X\subseteq S\).
Der rechte Nulltest fragt: Ist \(AX=\{0\}\)?
Der linke Nulltest fragt: Ist \(XA=\{0\}\)?
\emph{Den Test bestehen} bedeutet, dass die jeweilige Antwort ja ist:
Alle dabei gebildeten Einzelprodukte sind Null.
Es genügt nicht, dass die Produktmenge neben anderen Werten auch \(0\)
enthält. Die Seite gehört zum Test dazu, denn im Allgemeinen müssen
\(AX\) und \(XA\) nicht übereinstimmen.

\section*{Ein Beispiel, das uns durch den Beweis begleitet}

\begin{NPReadingExample}{Dreierbeispiel: Multiplikation und Nulltests}
Auf \(S=\{0,c,a\}\) sei \(a^2=c\); alle anderen Produkte seien \(0\).
Die vollständige Multiplikationstafel ist
\[
 \begin{array}{c|ccc}
 \cdot&0&c&a\\\hline
 0&0&0&0\\
 c&0&0&0\\
 a&0&0&c
 \end{array}
\]
Die Operation ist assoziativ: Bei drei Elementen \(x,y,z\) liegt sowohl
\(xy\) als auch \(yz\) in \(\{0,c\}\). Sowohl \(0\) als auch \(c\)
ergeben mit jedem weiteren Element auf beiden Seiten das Produkt \(0\).
Also sind \((xy)z\) und \(x(yz)\) beide gleich \(0\).

Die sieben nichtleeren Teilmengen lassen sich in zwei Familien schreiben:
\[
 \begin{aligned}
 p_0&=\bigl\{\{0\},\{c\},\{0,c\}\bigr\},\\
 p_1&=\bigl\{\{a\},\{0,a\},\{c,a\},\{0,c,a\}\bigr\}.
 \end{aligned}
\]
Die erste Familie besteht aus den Mengen ohne \(a\), die zweite aus
den Mengen mit \(a\). Für nichtleere Teilmengen gilt
\[
 AB=\{0\}\quad\Longleftrightarrow\quad
 a\notin A\ \text{oder}\ a\notin B.
\]
Denn das einzige Paar mit einem von \(0\) verschiedenen Produkt ist
\((a,a)\). Enthalten beide Mengen \(a\), kommt im Mengenprodukt also
\(c\) vor. Fehlt \(a\) in wenigstens einer Menge, ist jedes Einzelprodukt
\(0\). Wegen der Nichtleerheit der Mengen ist dann \(AB\) genau \(\{0\}\).

Eine Menge aus \(p_0\) ergibt deshalb mit \emph{jeder} nichtleeren
Menge ein Nullprodukt, links wie rechts. Eine Menge aus \(p_1\) ergibt
genau mit den Mengen aus \(p_0\) ein Nullprodukt. Innerhalb derselben
Familie kann man die Mengen durch solche Tests nicht unterscheiden.
Zwischen den Familien gelingt das sehr wohl: Der Test mit \(X=\{a\}\)
ergibt für \(\{0\}\) ein Nullprodukt, für \(\{a\}\) hingegen \(\{c\}\).
Die beiden Familien werden genau die zwei Profile dieses Beispiels sein.
\end{NPReadingExample}

\section*{Die Beweisidee am Beispiel}

\noindent\textbf{Behauptung.} Sind die Potenzhalbgruppen von \(S\) und
\(T\) isomorph und ist \(S\) endlich mit Null und \(S^2=\{0,c\}\),
\(c\ne0\), so sind auch \(S\) und \(T\) isomorph.

Was müssen wir dafür aus der großen Multiplikationstafel der Teilmengen
zurückgewinnen? In der kleinen Tafel des Dreierbeispiels gibt es zwei
Elemente, nämlich \(0,c\), die mit jedem Element auf beiden Seiten Null
ergeben. Das dritte Element \(a\) verhält sich anders: \(a^2\ne0\).
Wenn wir auf der Zielseite ebenfalls zwei Elemente der ersten Art und
eines der zweiten Art finden, können wir sie passend zuordnen.

\medskip\noindent\textbf{Was wird gezählt?}
Die Grundmenge \(S=\{0,c,a\}\) hat drei Elemente. Ihre
Potenzhalbgruppe hat dagegen sieben Elemente. Diese sieben Elemente
sind die folgenden \emph{Mengen}:
\[
 \begin{array}{c|l|c}
 \text{Größe der Teilmenge}&\text{Teilmengen von }S&\text{Anzahl}\\\hline
 1&\{0\},\ \{c\},\ \{a\}&3\\
 2&\{0,c\},\ \{0,a\},\ \{c,a\}&3\\
 3&\{0,c,a\}&1
 \end{array}
\]
Zusammen sind es \(3+3+1=7\) nichtleere Teilmengen.
Unter diesen sieben Mengen bestehen genau drei jeden linken und
rechten Nulltest: \(\{0\},\{c\},\{0,c\}\).
Sie sind sämtliche nichtleeren Teilmengen von \(\{0,c\}\).

\medskip\noindent\textbf{Wie kommt man von Mengenanzahlen zu Elementanzahlen?}
Für eine endliche Grundmenge mit \(n\) Elementen gibt es für jedes
Element zwei Entscheidungen: in die Teilmenge aufnehmen oder weglassen.
Das gibt \(2^n\) Teilmengen, darunter genau eine leere.
Also gibt es \(2^n-1\) nichtleere Teilmengen. Die ersten Werte sind
\[
 \begin{array}{c|ccccc}
 \text{Anzahl der Grundelemente }n&0&1&2&3&4\\\hline
 \text{Anzahl der nichtleeren Teilmengen }2^n-1&0&1&3&7&15
 \end{array}
\]
Diese Zahlen wachsen streng: Fügt man einer endlichen Grundmenge ein
Element hinzu, bleiben alle bisherigen Teilmengen erhalten, und es
kommen neue hinzu, etwa die Einermenge des neuen Elements.
Darum bestimmen drei nichtleere Teilmengen eine Zweiermenge und sieben
eine Dreiermenge, \emph{sofern jeweils alle nichtleeren Teilmengen}
der betreffenden Grundmenge gezählt werden.
Eine beliebige Auswahl von drei oder sieben Teilmengen würde diesen
Rückschluss nicht erlauben.

\medskip\noindent\textbf{Was bedeutet dann \(3-2=1\)?}
Jetzt zählen wir wieder Grundelemente: In \(S=\{0,c,a\}\) gibt es
insgesamt drei. Davon gehören zwei zu \(\{0,c\}\). Es bleibt genau
ein Grundelement außerhalb, nämlich \(a\).
Die sieben Teilmengen liefern die Grundelementzahl drei; die drei
stets nullmachenden Teilmengen liefern die Grundelementzahl zwei.
Mit diesen beiden Grundelementzahlen rechnen wir \(3-2=1\).

Damit dieselbe Rechnung auf der Zielseite funktioniert, sind noch
Beweise nötig. Zuerst gewinnen wir eine Null in \(T\).
Dann zeigen wir, dass \(\Phi\) Nulltests erhält und daher die drei
Mengen, die jeden Test bestehen, auf genau solche Mengen in \(T\)
abbildet. Schließlich begründen wir, dass diese Zielfamilie aus
\emph{allen} nichtleeren Teilmengen derjenigen Zielelemente besteht,
die mit jedem Element links und rechts Null ergeben.
Erst dann dürfen wir aus ihren drei Mitgliedern auf zwei solche
Zielelemente schließen. Alle sieben Teilmengen zusammen liefern
entsprechend drei Zielelemente insgesamt; eines bleibt außerhalb.

Allgemein kann es mehr als zwei Arten von Testverhalten geben.
Wir zählen dann zunächst die Grundelemente, die mindestens alle Tests
eines festen Verhaltens bestehen. Anschließend entfernen wir die
Elemente mit zusätzlichen bestandenen Tests. Die späteren Profile
und ihre Ordnung machen dieses Vorgehen präzise.

Aus den übereinstimmenden Anzahlen wollen wir schließlich eine
Bijektion \(f\colon S\to T\) der Grundelemente gewinnen.
Wir werden dafür zeigen, dass \(T\) eine Null \(0_T\) besitzt und
\(T^2=\{0_T,c_T\}\) mit \(c_T\ne0_T\) gilt. Die Bijektion soll
\[
 f(0)=0_T,\qquad f(c)=c_T
\]
erfüllen und für alle \(x,y\in S\) die Nullprodukte in beiden
Richtungen erhalten:
\[
 xy=0\quad\Longleftrightarrow\quad f(x)f(y)=0_T.
\]
Diese Aussagen sind an dieser Stelle noch zu beweisen.

\medskip\noindent\textbf{Warum würde das zum Ziel führen?}
Gesucht ist eine Bijektion, die außerdem
\(f(xy)=f(x)f(y)\) für jedes Paar \(x,y\) erfüllt.
Da \(xy\) nur die Werte \(0\) und \(c\) annehmen kann, genügen zwei Fälle.
Ist \(xy=0\), wären sowohl \(f(xy)=f(0)\) als auch \(f(x)f(y)\)
gleich \(0_T\). Ist \(xy=c\), müsste \(f(x)f(y)\ne0_T\) sein:
Andernfalls würde die Rückrichtung der angekündigten Äquivalenz
\(xy=0\) erzwingen. Als Produkt in \(T\) wäre \(f(x)f(y)\) dann
der einzige andere Produktwert \(c_T\), also
\[
 f(x)f(y)=c_T=f(c)=f(xy).
\]
Damit wäre in beiden Fällen die erforderliche Produktgleichheit
erfüllt. Zusammen mit der Bijektivität wäre \(f\) der gesuchte
Isomorphismus. Die folgenden Abschnitte beweisen die dafür
angekündigten Aussagen.

\section*{Zuerst klären wir, welche Eigenschaften \(T\) besitzt}

Bevor wir Nulltests auf beiden Seiten vergleichen, müssen wir wissen,
dass auch \(T\) eine Null besitzt. Dieser Teil des Beweises verwendet
zunächst nur die Produkte in den Potenzhalbgruppen.

Zunächst ist \(T\) nichtleer und endlich. Die Menge \(\PowNE S\)
ist nichtleer und endlich, also gilt wegen der Bijektivität von \(\Phi\)
dasselbe für \(\PowNE T\). Eine nichtleere Teilmenge von \(T\) kann
es nur geben, wenn \(T\) selbst nichtleer ist. Außerdem ordnet
\(t\mapsto\{t\}\) verschiedenen Elementen von \(T\) verschiedene
Elemente der endlichen Menge \(\PowNE T\) zu. Daher ist \(T\) endlich.

\par\Needspace{9\baselineskip}
\medskip\noindent\textbf{Nächstes Ziel: \(T^2\) hat genau zwei Elemente.}
Dazu zählen wir zunächst die verschiedenen \emph{Mengen}, die als
Ergebnis einer Mengenmultiplikation auftreten. Wir geben dieser
Familie einen eigenen Namen:
\[
 \begin{aligned}
 \mathcal M_S&:=\{AB\mid A,B\in\PowNE S\}=(\PowNE S)^2,\\
 \mathcal M_T&:=\{CD\mid C,D\in\PowNE T\}=(\PowNE T)^2.
 \end{aligned}
\]
Ein Mitglied von \(\mathcal M_S\) ist also eine Menge \(AB\).
Die Familie \(\PowNE S\) enthält dagegen alle erlaubten Eingaben
\(A\) und \(B\), auch solche, die niemals als Ergebnis vorkommen.
Im Dreierbeispiel ist etwa \(\{a\}\) eine erlaubte Eingabe, kann aber
kein Ergebnis sein, weil kein Einzelprodukt den Wert \(a\) hat.

\medskip\noindent\textbf{1. Auf der Ausgangsseite gibt es drei Produktmengen.}
Weil jedes Einzelprodukt in \(\{0,c\}\) liegt, ist jedes nichtleere
Mengenprodukt eine der Mengen \(\{0\},\{c\},\{0,c\}\).
Alle drei treten auf: Da \(c\in S^2\), gibt es \(u,v\in S\)
mit \(uv=c\), und dann gilt
\[
 \{0\}\{0\}=\{0\},\qquad
 \{u\}\{v\}=\{c\},\qquad
 \{0,u\}\{v\}=\{0,c\}.
\]
Damit haben wir genau die möglichen \emph{Ergebnisse} aufgezählt:
\[
 \mathcal M_S=\bigl\{\{0\},\{c\},\{0,c\}\bigr\}
             =\PowNE{(S^2)},\qquad |\mathcal M_S|=3.
\]
Im Dreierbeispiel hat \(\PowNE S\) sieben Mitglieder, die Familie
\(\mathcal M_S\) aber nur diese drei.

\medskip\noindent\textbf{2. Auf der Zielseite gibt es ebenfalls drei Produktmengen.}
Wir zeigen, dass die gegebene Abbildung \(\Phi\) eine Bijektion
zwischen den beiden Ergebnisfamilien liefert. Dazu betrachten wir
\[
 \Phi_2\colon\mathcal M_S\longrightarrow\mathcal M_T,
 \qquad U\longmapsto\Phi(U).
\]
Dies ist dieselbe Zuordnung wie \(\Phi\), nur auf die Produktmengen
beschränkt. Dafür sind drei Punkte zu prüfen.
\begin{itemize}
 \item \emph{Die Bilder liegen in \(\mathcal M_T\).}
       Ist \(U\in\mathcal M_S\), so ist \(U=AB\) für geeignete
       nichtleere \(A,B\). Dann ist
       \(\Phi(U)=\Phi(A)\Phi(B)\) ein Mengenprodukt in \(T\).
 \item \emph{Verschiedene Eingaben bleiben verschieden.}
       Aus \(\Phi_2(U)=\Phi_2(V)\) folgt
       \(\Phi(U)=\Phi(V)\), also \(U=V\), weil \(\Phi\) injektiv ist.
 \item \emph{Jede Menge aus \(\mathcal M_T\) wird getroffen.}
       Sei \(V=CD\in\mathcal M_T\). Die Surjektivität von \(\Phi\)
       liefert \(A,B\) mit \(\Phi(A)=C\) und \(\Phi(B)=D\).
       Dann liegt \(AB\) in \(\mathcal M_S\), und
       \(\Phi_2(AB)=\Phi(A)\Phi(B)=CD=V\).
\end{itemize}
Damit ist \(\Phi_2\) bijektiv, und es folgt
\(\lvert\mathcal M_T\rvert=\lvert\mathcal M_S\rvert=3\).
Die Produkterhaltung sorgt dabei für die passende Einschränkung;
aus der Bijektivität allein könnte man sie nicht folgern.

\par\Needspace{8\baselineskip}
\medskip\noindent\textbf{3. Drei Produktmengen erzwingen zwei Produktwerte.}
Jetzt schließen wir von der Anzahl der Mengen in \(\mathcal M_T\)
auf die Anzahl der Grundelemente in \(T^2\).
Jedes \(t\in T^2\) liefert ein Mitglied \(\{t\}\in\mathcal M_T\):
Es gibt \(u,v\in T\) mit \(uv=t\), also
\(\{u\}\{v\}=\{t\}\). Außerdem gehört \(T^2\) selbst zu
\(\mathcal M_T\), denn \(TT=T^2\) ist ebenfalls ein Mengenprodukt.
Hier dürfen wir \(T\) als Faktor wählen, weil \(T\) nichtleer ist.

Hätte \(T^2\) mindestens drei Elemente, könnten wir verschiedene
\(r,s,t\in T^2\) wählen. Dann hätte \(\mathcal M_T\) mindestens
die folgenden vier verschiedenen Mitglieder:
\[
 \underbrace{\boxed{\{r\}}\quad\boxed{\{s\}}\quad\boxed{\{t\}}}_{%
   \text{drei verschiedene Einermengen}}
 \qquad
 \underbrace{\boxed{T^2}}_{\text{enthält mindestens }r,s,t}.
\]
Die vierte Menge ist keine Einermenge und deshalb von den ersten
drei verschieden. Das widerspricht \(\lvert\mathcal M_T\rvert=3\).
Gezählt werden hier Mitglieder der Familie \(\mathcal M_T\):
Auch die gesamte Menge \(T^2\) ist genau ein solches Mitglied.
Also hat \(T^2\) höchstens zwei Elemente.

\(T^2\) ist nicht leer: Für ein \(u\in T\) gehört \(uu\) dazu.
Hätte \(T^2\) nur ein Element \(r\), wäre jedes Einzelprodukt
gleich \(r\). Jedes nichtleere Mengenprodukt wäre dann genau
\(\{r\}\). Das ergäbe nur ein Mitglied von \(\mathcal M_T\),
wieder im Widerspruch zu deren Größe drei. Somit gilt
\[
 \boxed{|T^2|=2.}
\]

\medskip\noindent\textbf{4. Die drei Produktmengen sind jetzt ausdrücklich bekannt.}
Schreiben wir \(T^2=\{r,s\}\) mit \(r\ne s\).
Schritt~3 hat bereits gezeigt, dass \(\{r\}\), \(\{s\}\) und
\(T^2=\{r,s\}\) zu \(\mathcal M_T\) gehören.
Weitere Mitglieder gibt es nach Schritt~2 nicht. Also ist
\[
 \mathcal M_T=\bigl\{\{r\},\{s\},\{r,s\}\bigr\}
             =\PowNE{(T^2)}.
\]
Diese letzte Gleichheit ist eine \emph{Folge} der bewiesenen Zweiergröße.
Wir brauchen sie im nächsten Abschnitt, um die Multiplikation
innerhalb der beiden Produktquadrate zu vergleichen.
Der eben geführte Beweis ist in
\FormulaRefAuto{NPTwoProductValuesTransport}[thm] vollständig tabelliert.

\section*{Wie wir aus dem Zweierquadrat eine Null gewinnen}

Wir wissen jetzt, dass \(T^2\) zwei Elemente hat. Noch ist keines
davon als Null ausgezeichnet. Zunächst suchen wir eine Null
\emph{innerhalb} von \(T^2\); danach zeigen wir, dass sie auch
gegenüber allen übrigen Elementen von \(T\) absorbierend wirkt.
Die Zweiergröße allein liefert noch keine Null. Wir verwenden jetzt
zusätzlich die Multiplikation auf \(S^2\) und den eingeschränkten
Isomorphismus der Potenzhalbgruppen.

Die Menge \(S^2=\{0,c\}\) ist unter Multiplikation abgeschlossen,
ebenso \(T^2\): Das Produkt zweier Elemente des Quadrats ist insbesondere
ein Produkt zweier Elemente der ursprünglichen Halbgruppe und liegt
deshalb wieder im Quadrat. Wir können beide Quadrate als Halbgruppen
betrachten. Nach den beiden Gleichheiten des vorigen Abschnitts
ist die dort bereits geprüfte Bijektion \(\Phi_2\) eine Abbildung
\[
 \Phi_2\colon\PowNE{(S^2)}\longrightarrow\PowNE{(T^2)}.
\]
Auch die Produkterhaltung bleibt erhalten: Für nichtleere
\(A,B\subseteq S^2\) ist \(AB\) wieder eine nichtleere Teilmenge
von \(S^2\), und
\(\Phi_2(AB)=\Phi(AB)=\Phi(A)\Phi(B)=\Phi_2(A)\Phi_2(B)\).
Somit ist \(\Phi_2\) ein Isomorphismus dieser kleineren Potenzhalbgruppen.
Im Quadrat \(S^2\) sind drei der vier Produkte schon bekannt:
\(00=0c=c0=0\). Nur \(c^2\) ist noch offen. Weil auch dieser Wert
in \(S^2\) liegt, gibt es genau die folgenden zwei Fälle.

\medskip\noindent\textbf{Erster Fall: \(c^2=0\).}
Dann sind alle Produkte zweier Elemente von \(S^2\) gleich \(0\).
Also ist jedes Produkt zweier nichtleerer Teilmengen von \(S^2\)
dieselbe Menge \(\{0\}\). Seien \(U,V\) nichtleere Teilmengen von
\(T^2\). Weil \(\Phi_2\) surjektiv ist, gibt es nichtleere
\(A,B\subseteq S^2\) mit \(U=\Phi_2(A)\) und \(V=\Phi_2(B)\).
Also ist \(UV=\Phi_2(AB)=\Phi_2(\{0\})\).
Alle diese Produkte haben denselben Wert, den wir \(Z\) nennen.

Wählen wir ein Element \(u\in T^2\), so ist insbesondere
\(Z=\{u\}\{u\}=\{u^2\}\). \(Z\) ist also eine Einermenge.
Wir nennen ihr einziges Element \(z\). Aus \(z=u^2\) folgt auch
\(z\in T^2\), denn \(u^2\) ist ein Produkt zweier Elemente von \(T\).
Für beliebige \(v,w\in T^2\)
gilt nun \(\{vw\}=\{v\}\{w\}=Z=\{z\}\), folglich \(vw=z\).
Setzen wir \(v=z\) oder
\(w=z\), erhalten wir \(zw=wz=z\) für jedes \(w\in T^2\).
Das heißt genau: \(z\) ist eine Null von \(T^2\).

\medskip\noindent\textbf{Zweiter Fall: \(c^2=c\).}
Hier lautet die Tafel der nichtleeren Teilmengen von \(S^2\)
\[
 \begin{array}{c|ccc}
 \cdot&\{0\}&\{c\}&\{0,c\}\\\hline
 \{0\}&\{0\}&\{0\}&\{0\}\\
 \{c\}&\{0\}&\{c\}&\{0,c\}\\
 \{0,c\}&\{0\}&\{0,c\}&\{0,c\}
 \end{array}
\]
Die Tafel ist symmetrisch, also ist die Multiplikation kommutativ.
Auf der Diagonale steht jeweils die Eingabe selbst: Jede der drei
\emph{Teilmengen} \(A\) erfüllt \(AA=A\), ist also ein idempotentes
Element der Potenzhalbgruppe von \(S^2\).
Für beliebige nichtleere \(U,V\subseteq T^2\) wählen
wir wieder Urbilder \(A,B\) unter \(\Phi_2\). Dann gilt
\[
 \begin{aligned}
 UV&=\Phi_2(AB)=\Phi_2(BA)=VU,\\
 UU&=\Phi_2(AA)=\Phi_2(A)=U.
 \end{aligned}
\]
Damit sind Kommutativität und Idempotenz auf der Zielseite bewiesen.

Nun gehen wir von diesen Aussagen über Teilmengen zu Aussagen über
Grundelemente über, indem wir Einermengen einsetzen. Aus
\(\{u\}\{v\}=\{v\}\{u\}\) folgt \(uv=vu\), und aus
\(\{u\}\{u\}=\{u\}\) folgt \(u^2=u\). Schreiben wir
\(T^2=\{u,v\}\), so setzen wir \(z=uv\).
Als Produkt zweier Elemente von \(T\) gehört \(z\) zu \(T^2\).
Es ist eine Null dieses Quadrats, denn
\[
 \begin{aligned}
 uz&=u(uv)=(u^2)v=uv=z,\\
 vz&=v(uv)=u(v^2)=uv=z.
 \end{aligned}
\]
Die erste Zeile benutzt Assoziativität und \(u^2=u\). In der zweiten
Zeile dürfen wir wegen der Kommutativität die Faktoren umordnen;
anschließend verwenden wir \(v^2=v\). Weil \(u,v\) alle Elemente von
\(T^2\) sind, haben wir damit die Multiplikation mit jedem Element des
Quadrats geprüft. Die Produkte \(zu,zv\) sind wegen der Kommutativität
ebenfalls gleich \(z\).

\medskip\noindent\textbf{Warum genügt eine Null des Quadrats?}
Beide Fälle haben dieselbe, zunächst auf das Quadrat beschränkte
Aussage geliefert: Es gibt \(z\in T^2\) mit \(zw=wz=z\) für
jedes \(w\in T^2\). Weil \(z\) selbst zu \(T^2\) gehört, dürfen
wir hier \(w=z\) einsetzen und erhalten \(zz=z\).

Nun wählen wir ein beliebiges \(t\in T\). Es könnte außerhalb von
\(T^2\) liegen; deshalb dürfen wir die bisherige Aussage nicht einfach
mit \(w=t\) anwenden. Wir verwenden stattdessen das Element \(w=zt\).
Dieses liegt sicher in \(T^2\): Es ist das Produkt der beiden
Elemente \(z,t\in T\). Daher wissen wir bereits \(z(zt)=z\).
Durch Assoziativität und \(zz=z\) ist \(z(zt)=(zz)t=zt\).
Zusammen folgt \(zt=z\). Für \(tz\) funktioniert derselbe Schritt
mit \(w=tz\in T^2\). In Rechenzeilen:
\[
 \begin{aligned}
  zt&=(zz)t&&\text{weil }zz=z\\
    &=z(zt)&&\text{durch Assoziativität}\\
    &=z&&\text{weil }zt\in T^2,\\[4pt]
  tz&=t(zz)&&\text{weil }zz=z\\
    &=(tz)z&&\text{durch Assoziativität}\\
    &=z&&\text{weil }tz\in T^2.
 \end{aligned}
\]
Wir haben also nicht vorausgesetzt, dass \(zt\) oder \(tz\) schon
gleich \(z\) ist. Es genügte zunächst ihre Zugehörigkeit zu \(T^2\);
die Gleichheiten folgten erst aus der Rechnung. Weil \(t\) beliebig
war, ist \(z\) eine Null von ganz \(T\). Dies ist genau
\FormulaRefAuto{NPZeroFromSquareAbsorption}[thm].

Wir schreiben von jetzt an \(0_T=z\) und nennen das andere Element
von \(T^2\) \(c_T\). Der vorbereitende Teil ist damit abgeschlossen:
Beide Halbgruppen sind endlich, beide besitzen eine Null, und es gilt
\[
 S^2=\{0,c\},\qquad T^2=\{0_T,c_T\},\qquad c_T\ne0_T.
\]
Die gesamte Herleitung der Zielnull steht in
\FormulaRefAuto{NPTargetZero}[thm]:
\NullproductsProofLink{NPTargetZero}.

\section*{Nulltests vergleichen, bevor wir Mengen zusammenfassen}

Wir verwenden nun die vor dem Dreierbeispiel eingeführten Nulltests
systematisch: \(A\) besteht den rechten Test mit \(X\), wenn
\(AX=\{0\}\) ist, und den linken, wenn \(XA=\{0\}\) ist.
Verglichen werden jeweils die Antworten bei \emph{derselben}
Testmenge und auf \emph{derselben} Seite.

\begin{NPReadingExample}{Dreierbeispiel: Einen Nulltest ausrechnen}
Wir testen \(A=\{0,a\}\) mit zwei verschiedenen rechten Testmengen:
\[
 \begin{aligned}
 A\{a\}&=\{0a,aa\}=\{0,c\}&&\text{Test nicht bestanden},\\
 A\{c\}&=\{0c,ac\}=\{0\}&&\text{Test bestanden}.
 \end{aligned}
\]
In der ersten Zeile kommt zwar \(0\) vor, aber auch \(c\ne0\).
Erst in der zweiten Zeile ist wirklich jedes Einzelprodukt Null.
\end{NPReadingExample}

Zunächst wollen wir sagen können, dass \(A\) mindestens alle Nulltests
von \(B\) besteht. Dafür schreiben wir \(A\preceq_S B\) und definieren
\[
 \begin{aligned}
 A\preceq_S B\quad\Longleftrightarrow\quad
 &\forall X\in\PowNE S\;\Bigl(
 (BX=\{0\}\Rightarrow AX=\{0\})\\[-2pt]
 &\hspace{7em}\land(XB=\{0\}\Rightarrow XA=\{0\})\Bigr).
 \end{aligned}
\]
\(A\preceq_S B\) heißt also: Jeder von \(B\) bestandene Test wird
auch von \(A\) bestanden; \(A\) darf zusätzliche Tests bestehen.
Dies ist eine Aussage über jeden einzelnen Test, keine bloße
Ungleichung zwischen Anzahlen bestandener Tests.
Insbesondere gilt \(\{0\}\preceq_S B\) für
jede nichtleere Menge \(B\), denn \(\{0\}\) besteht alle Nulltests.
Das Zeichen \(\preceq_S\) ist keine Behauptung über die Inklusion
von \(A\) in \(B\).

Dieser Vergleich ist reflexiv: Besteht \(A\) einen Test, so besteht
\(A\) denselben Test. Er ist auch transitiv. Aus \(A\preceq_S B\)
und \(B\preceq_S C\) folgt für jeden rechten Nulltest
\[
 CX=\{0\}\ \Longrightarrow\ BX=\{0\}\ \Longrightarrow\ AX=\{0\}.
\]
Für linke Nulltests gilt dieselbe Schlussfolge. Also ist
\(A\preceq_S C\). Eine reflexive und transitive Relation heißt eine
Präordnung. Sie muss verschiedene Mengen noch nicht unterscheiden:
Im Beispiel bestehen \(\{0\}\) und \(\{c\}\) genau dieselben Tests,
obwohl die Mengen verschieden sind.

Wir schreiben \(A\sim_S B\), wenn beide Vergleichsrichtungen gelten,
also wenn \(A\) und \(B\) genau dieselben rechten und linken Nulltests
bestehen. Ausgeschrieben bedeutet das
\[
 \begin{aligned}
 A\sim_S B\quad\Longleftrightarrow\quad
 &\forall X\in\PowNE S\;\Bigl(
 (AX=\{0\}\Longleftrightarrow BX=\{0\})\\[-2pt]
 &\hspace{7em}\land(XA=\{0\}\Longleftrightarrow XB=\{0\})\Bigr).
 \end{aligned}
\]
Diese Relation ist eine Äquivalenzrelation. Reflexivität und Symmetrie
sind unmittelbar in der Formulierung enthalten. Für die Transitivität
beachten wir: Stimmen die Antworten von \(A\) und \(B\) bei jedem Test
überein und ebenso die Antworten von \(B\) und \(C\), so stimmen auch
die Antworten von \(A\) und \(C\) überein.

\section*{Der Bezug zu den Green-Relationen}

Dieser Abschnitt ordnet die Nulltests in bekannte Begriffe ein;
der anschließende Beweis verwendet nur die bereits definierten Tests.
Die Frage nach den Green-Relationen aus Band~28 liegt nahe: Auch dort
werden Elemente einer Halbgruppe mithilfe ihres Multiplikationsverhaltens
verglichen und in Klassen eingeteilt. Wir schreiben die Definitionen
hier noch einmal aus.

Sei \(H\) eine Halbgruppe, deren Multiplikation wir durch
Nebeneinandersetzen schreiben, und \(u\in H\). Die von \(u\)
erzeugten linken, rechten und zweiseitigen Hauptideale sind
\[
 \begin{aligned}
 L_H(u)&=\{u\}\cup\{xu\mid x\in H\},\\
 R_H(u)&=\{u\}\cup\{uy\mid y\in H\},\\
 J_H(u)&=L_H(u)\cup R_H(u)\cup\{xuy\mid x,y\in H\}.
 \end{aligned}
\]
Man nimmt also den Erzeuger selbst sowie die durch Multiplikation
auf der jeweils zugelassenen Seite entstehenden Elemente.
\(u\) wird ausdrücklich aufgenommen, weil \(H\) kein neutrales
Element besitzen muss. Dies sind die Definitionen
\FormulaRefAuto{SemigroupPrincipalLeftIdealDef}[def],
\FormulaRefAuto{SemigroupPrincipalRightIdealDef}[def] und
\FormulaRefAuto{SemigroupPrincipalTwoSidedIdealDef}[def] aus Band~28,
mit unterdrücktem Index für die Multiplikation.

Für \(u,v\in H\) sind nun die Green-Relationen definiert durch
\[
 \begin{aligned}
 u\mathrel{\mathcal L_H}v&\quad\Longleftrightarrow\quad L_H(u)=L_H(v),\\
 u\mathrel{\mathcal R_H}v&\quad\Longleftrightarrow\quad R_H(u)=R_H(v),\\
 u\mathrel{\mathcal J_H}v&\quad\Longleftrightarrow\quad J_H(u)=J_H(v),\\
 u\mathrel{\mathcal H_H}v&\quad\Longleftrightarrow\quad
   u\mathrel{\mathcal L_H}v\ \text{und}\ u\mathrel{\mathcal R_H}v.
 \end{aligned}
\]
Verglichen werden also die erzeugten Hauptideale; für
\(\mathcal H_H\) müssen sowohl die linken als auch die rechten
Hauptideale übereinstimmen. Die zugehörigen Definitionen sind
\FormulaRefAuto{SemigroupGreenLDef}[def],
\FormulaRefAuto{SemigroupGreenRDef}[def],
\FormulaRefAuto{SemigroupGreenJDef}[def] und
\FormulaRefAuto{SemigroupGreenHDef}[def].
Für den Vergleich mit unseren Profilen ist \(H\) die
\emph{Potenzhalbgruppe} von \(S\): Ihre Elemente \(u,v\) sind dann
nichtleere Teilmengen von \(S\), und die Produkte sind Mengenprodukte.

Bei unseren Nulltests fragen wir dagegen, welche Elemente eine Menge
auf der jeweiligen Seite zu Null machen. Dies ist unmittelbar der
Begriff des \emph{Annihilators} aus Band~33:
\[
 \begin{aligned}
 \operatorname{Ann}_r(A)&=\{x\in S\mid ax=0\text{ für alle }a\in A\},\\
 \operatorname{Ann}_l(A)&=\{x\in S\mid xa=0\text{ für alle }a\in A\}.
 \end{aligned}
\]
Vgl. \FormulaRefAuto{SemigroupRightAnnihilatorDef}[def] und
\FormulaRefAuto{SemigroupLeftAnnihilatorDef}[def].
Für nichtleere \(A,X\) gilt
\[
 AX=\{0\}\ \Longleftrightarrow\ X\subseteq\operatorname{Ann}_r(A),
 \qquad
 XA=\{0\}\ \Longleftrightarrow\ X\subseteq\operatorname{Ann}_l(A).
\]
Denn beide Seiten besagen jeweils, dass alle betreffenden
Einzelprodukte Null sind; die Nichtleerheit sichert, dass die
Produktmenge tatsächlich \(\{0\}\) ist. Deshalb gilt genau
\[
 A\preceq_S B\quad\Longleftrightarrow\quad
 \begin{cases}
 \operatorname{Ann}_r(B)\subseteq\operatorname{Ann}_r(A),\\
 \operatorname{Ann}_l(B)\subseteq\operatorname{Ann}_l(A).
 \end{cases}
\]
Von links nach rechts setzt man die einzelnen Tests \(X=\{x\}\)
ein. Von rechts nach links benutzt man die gerade gezeigten
Teilmengenbedingungen. Zwei Mengen haben also genau dann dasselbe
Nullprofil, wenn beide Annihilatoren übereinstimmen.
Die Richtung der Inklusion erklärt auch unser Ordnungszeichen:
\(A\preceq_S B\) bedeutet, dass jedes Element, das \(B\) auf einer
Seite zu Null macht, auch \(A\) auf derselben Seite zu Null macht.
Es geht um die Zugehörigkeit der einzelnen Elemente, nicht nur um
die Anzahl der nullmachenden Elemente.

\begin{NPReadingExample}{Dreierbeispiel: Nullprofile und Green-Klassen}
Die Mengen \(\{0\}\) und \(\{c\}\) bestehen hier dieselben Nulltests:
Beide haben links und rechts ganz \(S\) als Annihilator. Sie besitzen
also dasselbe Nullprofil. Als Elemente der \emph{Potenzhalbgruppe}
erzeugen sie jedoch verschiedene zweiseitige Hauptideale.
Das von \(\{0\}\) erzeugte Hauptideal ist \(\{\{0\}\}\).
Das von \(\{c\}\) erzeugte Hauptideal ist
\(\{\{c\},\{0\}\}\): Der Erzeuger selbst gehört dazu, und jede
Multiplikation mit einem weiteren Faktor ergibt \(\{0\}\).
Die beiden Elemente sind daher nicht Green-\(\mathcal J\)-verwandt.
Unsere Profilgleichheit ist folglich keine Umbenennung der
Green-\(\mathcal J\)-Relation. Dasselbe Beispiel trennt sie auch
von \(\mathcal L\), \(\mathcal R\) und \(\mathcal H\), denn
bereits die einseitigen Hauptideale sind verschieden.
\end{NPReadingExample}
Für den weiteren Beweis benötigen wir nur die Nulltests und ihre
Präordnung. Die Green-Relationen müssen wir dazu nicht weiter untersuchen.

\par\Needspace{14\baselineskip}
\section*{Was \([A]_S\) und \(Q_S\) bedeuten}

Nun fassen wir die Mengen mit denselben Testergebnissen zusammen:
\[
 [A]_S=\{B\in\PowNE S\mid B\sim_S A\}.
\]
Dies ist die \emph{Äquivalenzklasse} von \(A\): die Familie aller
nichtleeren Teilmengen von \(S\), die dieselben Nulltests wie \(A\)
bestehen. Wir nennen diese Klasse das Nullproduktprofil von \(A\).
\(A\) ist eine einzelne Teilmenge; \([A]_S\) ist eine Familie solcher
Teilmengen. Sprechen wir vom \emph{Profil eines Grundelements}
\(s\in S\), meinen wir stets \([\{s\}]_S\), also das Profil seiner
Einermenge. Das Grundelement \(s\), die Menge \(\{s\}\) und deren
Profil haben damit jeweils eine eigene Rolle.

\par\Needspace{14\baselineskip}
\begin{NPReadingExample}{Dreierbeispiel: Die beiden Profilklassen}
Im Dreierbeispiel ist
\[
 [\{0\}]_S=[\{c\}]_S=[\{0,c\}]_S=p_0,
\]
denn diese drei Mengen bestehen jeden Nulltest. Die vier Mengen mit
\(a\) besitzen dagegen alle das Profil \(p_1=[\{a\}]_S\).
Damit ist
\[
 Q_S=\{[A]_S\mid A\in\PowNE S\}=\{p_0,p_1\}
 \qquad\text{im Beispiel}.
\]
\end{NPReadingExample}
Allgemein bezeichnet \(Q_S\) die Menge \emph{aller} Nullproduktprofile.
Die ebenfalls gebräuchliche Schreibweise
\(Q_S=(\PowNE S)/{\sim_S}\) nennt man einen Quotienten.

Auf den Profilen setzen wir
\[
 [A]_S\le[B]_S\quad\Longleftrightarrow\quad A\preceq_S B.
\]
\medskip\noindent\textbf{Warum hängt der Vergleich nur von den Klassen ab?}
Dieselbe Klasse kann durch verschiedene Mengen beschrieben werden.
Seien \([A']_S=[A]_S\) und \([B']_S=[B]_S\). Dann bestehen \(A'\)
und \(A\) dieselben Tests, ebenso \(B'\) und \(B\).
Wir müssen zeigen, dass \(A\preceq_S B\) und \(A'\preceq_S B'\)
denselben Wahrheitswert haben.

Nehmen wir zunächst \(A\preceq_S B\) an und wählen eine beliebige
nichtleere Testmenge \(X\). Für einen von \(B'\) bestandenen rechten
Test gilt die Schlussfolge
\[
 B'X=\{0\}\ \Longrightarrow\ BX=\{0\}
 \ \Longrightarrow\ AX=\{0\}\ \Longrightarrow\ A'X=\{0\}.
\]
Der erste Schritt benutzt die gleichen Tests von \(B'\) und \(B\),
der zweite die Annahme \(A\preceq_S B\), der dritte die gleichen
Tests von \(A\) und \(A'\). Auf der linken Seite gilt entsprechend
\[
 XB'=\{0\}\ \Longrightarrow\ XB=\{0\}
 \ \Longrightarrow\ XA=\{0\}\ \Longrightarrow\ XA'=\{0\}.
\]
Somit besteht \(A'\) jeden von \(B'\) bestandenen Test, also
\(A'\preceq_S B'\). Für die Rückrichtung vertauschen wir die Rollen
von \(A,A'\) und von \(B,B'\); die gleichen Testantworten gelten
jeweils in beiden Richtungen. Der Vergleich ist deshalb unabhängig
von den gewählten Vertretern. Genau dies bedeutet hier
\emph{wohldefiniert}.

\medskip\noindent\textbf{Warum ist dies eine Halbordnung?}
Wir prüfen die drei dafür erforderlichen Eigenschaften auf den Klassen.
\begin{itemize}
 \item \emph{Reflexivität.} Sei \(p=[A]_S\) ein beliebiges Profil.
       Weil \(A\preceq_S A\) gilt, liefert die Definition
       \([A]_S\le[A]_S\), also \(p\le p\).
 \item \emph{Transitivität.} Seien \(p=[A]_S\), \(q=[B]_S\) und
       \(r=[C]_S\) mit \(p\le q\) und \(q\le r\).
       Wegen der gerade bewiesenen Vertreterunabhängigkeit dürfen wir
       in beiden Vergleichen dieselbe Menge \(B\) für \(q\) verwenden.
       Damit gelten \(A\preceq_S B\) und \(B\preceq_S C\).
       Die Transitivität der Präordnung liefert \(A\preceq_S C\),
       also \(p\le r\).
 \item \emph{Antisymmetrie.} Gelten \([A]_S\le[B]_S\) und
       \([B]_S\le[A]_S\), so haben wir \(A\preceq_S B\) und
       \(B\preceq_S A\). Daher bestehen \(A\) und \(B\) genau
       dieselben Tests, also \(A\sim_S B\).
       Jede Menge hat dann genau dann dieselben Tests wie \(A\),
       wenn sie dieselben Tests wie \(B\) hat. Folglich enthalten
       \([A]_S\) und \([B]_S\) dieselben Mengen und sind gleich.
\end{itemize}
Damit ist \(Q_S\) eine halbgeordnete Menge. Auf den einzelnen Mengen
war der Vergleich nur eine Präordnung, weil verschiedene Mengen
dieselben Tests bestehen können. Auf den Klassen werden gerade
solche Mengen zu einem einzigen Profil zusammengefasst.
Die allgemeinen Nachweise dieser Quotientenkonstruktion stehen in
\FormulaRefAuto{PreorderKernelCompatibility}[thm] und
\FormulaRefAuto{PreorderQuotientPartialOrder}[thm] aus Band~12.

\begin{NPReadingExample}{Dreierbeispiel: Die Profilordnung}
Im Beispiel gilt \(p_0<p_1\). Denn jede Menge aus \(p_0\) besteht alle
Tests, insbesondere diejenigen einer Menge aus \(p_1\).
Für die umgekehrte Richtung müsste \(\{a\}\preceq_S\{0\}\) gelten.
Der rechte Test mit \(X=\{a\}\) widerlegt das:
\(\{0\}X=\{0\}\), aber \(\{a\}X=\{c\}\ne\{0\}\).
\end{NPReadingExample}
Wir verwenden \(q<p\) fortan als Abkürzung für \(q\le p\) und \(q\ne p\).
Zwischen anderen Profilen muss eine Vergleichbarkeit keineswegs bestehen.

Dieselben Definitionen treffen wir für \(T\), mit \(0_T\) anstelle
von \(0\). So erhalten wir \(\preceq_T\), \([B]_T\) und \(Q_T\).
Der Index zeigt jeweils, in welcher Halbgruppe die Nulltests stattfinden.
Auf den Profilen führen wir nur eine Ordnung ein; eine Multiplikation
von Profilen wird für den Beweis nicht benötigt.

\section*{Drei Mengen mit verschiedenen Aufgaben}

Ein Profil \(p\in Q_S\) ist nun fest gewählt. Wir wollen zählen,
wie viele \emph{Grundelemente} genau dieses Profil besitzen.
Im Dreierbeispiel ist das für \(p_1\) unmittelbar sichtbar:
\[
 p_1=\bigl\{\{a\},\{0,a\},\{c,a\},\{0,c,a\}\bigr\}.
\]
Diese Familie hat vier Mitglieder, aber nur eines ist eine Einermenge,
nämlich \(\{a\}\). Deshalb besitzt genau ein Grundelement das Profil
\(p_1\). Auf der Zielseite können wir die Einermengen unter den
Bildmengen jedoch nicht einfach erkennen: \(\Phi\) muss Einermengen
nicht auf Einermengen abbilden. Die unmittelbar zugängliche Anzahl
der Mengen im Profil reicht deshalb noch nicht für unsere Zählaufgabe.

Wir suchen eine Familie, aus deren Größe sich die Zahl ihrer
Grundelemente mit \(2^n-1\) zurückgewinnen lässt. Dafür muss die
Familie \emph{alle} nichtleeren Teilmengen einer Grundmenge enthalten.
\(p_1\) erfüllt das nicht: Es enthält \(\{0,a\}\), aber nicht
\(\{0\}\). Wäre \(p_1\) die nichtleere Potenzmenge einer Menge
\(D\), müsste mit \(\{0,a\}\subseteq D\) auch \(0\in D\) gelten;
dann müsste \(\{0\}\) ebenfalls in \(p_1\) liegen.

\medskip\noindent\textbf{Warum nehmen wir gerade kleinere Profile hinzu?}
Verkleinern wir eine nichtleere Menge \(A\) auf eine nichtleere
Teilmenge \(B\), so kann kein von \(A\) bestandener Nulltest
verloren gehen: Sind alle Produkte mit Elementen aus \(A\) Null,
gilt das insbesondere für diejenigen mit Elementen aus \(B\).
Ein zuvor nicht bestandener Test kann dagegen hinzukommen.
In unserer Ordnung heißt das \(B\preceq_S A\), also
\([B]_S\le[A]_S\). Die beim Verkleinern fehlenden Teilmengen
gehören somit zum selben oder zu einem kleineren Profil.

Im Beispiel fehlen in \(p_1\) genau \(\{0\},\{c\},\{0,c\}\).
Diese bilden das kleinere Profil \(p_0\). Nehmen wir es hinzu, gilt
\[
 p_0\cup p_1=\PowNE S.
\]
Jetzt dürfen wir zurückrechnen: Die sieben Mengen dieser Familie
liefern drei Grundelemente insgesamt. Das kleinere Profil \(p_0\)
ist seinerseits die Familie aller drei nichtleeren Teilmengen von
\(\{0,c\}\) und liefert zwei Grundelemente.
Nach Abzug dieser beiden bleibt \(3-2=1\) Grundelement mit
\emph{genau} dem Profil \(p_1\). Gezählt werden also zuerst die
Grundelemente mit Profil \emph{kleiner oder gleich} \(p_1\);
danach werden diejenigen mit \emph{kleinerem} Profil abgezogen.

So wollen wir auch bei einer beliebigen Profilordnung vorgehen.
Wir sammeln zunächst sämtliche Mengen aus dem festen Profil und
allen kleineren Profilen. Dass diese Familie tatsächlich aus
\emph{allen} nichtleeren Teilmengen einer bestimmten Grundmenge
besteht, beweisen wir im nächsten Abschnitt. Erst danach dürfen
wir ihre Größe zur Rückrechnung verwenden.
Die folgenden drei Definitionen geben der gesammelten Familie,
ihrer Grundmenge und den Grundelementen mit genau einem Profil
jeweils einen eigenen Namen:
\[
 \begin{aligned}
 L_p^S&=\{A\in\PowNE S\mid[A]_S\le p\},\\
 D_p^S&=\{s\in S\mid[\{s\}]_S\le p\},\\
 C_p^S&=\{s\in S\mid[\{s\}]_S=p\}.
 \end{aligned}
\]
\(L_p^S\) sammelt \emph{nichtleere Teilmengen}, deren Profile unterhalb
von oder gleich \(p\) liegen. Dagegen sammelt \(D_p^S\)
\emph{Grundelemente}, deren Einermengen ein solches Profil besitzen.
\(C_p^S\) sammelt schließlich nur die Grundelemente mit Profil genau
\(p\). Ein \(s\) gehört also zu \(C_p^S\), wenn seine Einermenge
\(\{s\}\) als Mitglied in der Familie \(p\) vorkommt.

Insbesondere gilt \(L_p^S\subseteq\PowNE S\), nicht
\(L_p^S\subseteq Q_S\). Elemente von \(L_p^S\) sind Mengen \(A\),
Elemente von \(Q_S\) sind Klassen \([A]_S\). Die Menge aller Profile
\(q\le p\) wäre eine andere Menge. Die Hochstellung \(S\) ist hier
nur eine Kennzeichnung der Ausgangshalbgruppe, keine Potenz.
Für \(T\) definieren wir entsprechend \(L_r^T,D_r^T,C_r^T\),
jeweils mit \(r\in Q_T\), Grundelementen aus \(T\) und deren
Nulltests bezüglich \(0_T\).

\begin{NPReadingExample}{Dreierbeispiel: Teilmengen und Grundelemente zählen}
Im Dreierbeispiel können wir alles vollständig ausrechnen:
\[
 \begin{array}{c|c|c|c}
 p&L_p^S&D_p^S&C_p^S\\\hline
 p_0&p_0&\{0,c\}&\{0,c\}\\
 p_1&p_0\cup p_1&S&\{a\}
 \end{array}
\]
Unterhalb von oder gleich \(p_0\) liegt nur \(p_0\) selbst. Daher enthält
\(L_{p_0}^S\) die drei Mengen \(\{0\},\{c\},\{0,c\}\).
Genau die Einermengen von \(0\) und \(c\) haben dieses Profil,
sodass \(D_{p_0}^S=C_{p_0}^S=\{0,c\}\).

Unterhalb von oder gleich \(p_1\) liegen dagegen beide Profile.
Deshalb enthält \(L_{p_1}^S\) \emph{alle sieben} nichtleeren Teilmengen
und \(D_{p_1}^S\) alle drei Grundelemente. Unter den Einermengen
hat aber nur \(\{a\}\) das Profil \(p_1\). Darum ist \(C_{p_1}^S=\{a\}\).
Hier sind drei verschiedene Anzahlen zu unterscheiden:
\[
 |p_1|=4,\qquad |L_{p_1}^S|=7,\qquad |C_{p_1}^S|=1.
\]
Die vier Mengen im Profil sind nicht vier Grundelemente mit diesem
Profil. Diese Unterscheidung ist für den folgenden Beweis entscheidend.
\end{NPReadingExample}

\section*{Die nichtleeren Teilmengen von \(D_p^S\)}

Der entscheidende Zusammenhang lautet
\[
 \boxed{L_p^S=\PowNE{(D_p^S)}.}
\]
Er sagt: Eine nichtleere Menge hat genau dann ein Profil unterhalb
von oder gleich \(p\), wenn \emph{jedes ihrer Elemente} zu \(D_p^S\)
gehört. Wir beweisen beide Richtungen und erklären dabei, weshalb
aus den Tests einer Menge die Tests ihrer einzelnen Elemente hervorgehen.

Wählen wir zunächst einen Vertreter \(B\) des Profils \(p\), also
eine nichtleere Menge mit \([B]_S=p\). Einen solchen Vertreter gibt es,
weil jedes Profil per Definition die Klasse irgendeiner Menge ist.

\medskip\noindent\textbf{Von \(L_p^S\) zu den Teilmengen von \(D_p^S\).}
Sei \(A\in L_p^S\). Dann ist \(A\preceq_S B\).
Wir wählen ein beliebiges \(a\in A\) und zeigen \(\{a\}\preceq_S B\).
Sei dazu \(X\) eine nichtleere Teilmenge mit \(BX=\{0\}\).
Aus \(A\preceq_S B\) folgt \(AX=\{0\}\). Somit ist \(ax=0\)
für jedes \(x\in X\), denn jedes dieser Produkte kommt in \(AX\) vor.
Wegen \(X\ne\varnothing\) enthält \(\{a\}X\) tatsächlich einen
Produktwert. Daher ist \(\{a\}X=\{0\}\).

Für einen linken Test mit \(XB=\{0\}\) erhalten wir entsprechend
\(XA=\{0\}\), also \(xa=0\) für jedes \(x\in X\), und folglich
\(X\{a\}=\{0\}\). Damit sind beide erforderlichen Implikationen
gezeigt. Also ist \([\{a\}]_S\le p\), das heißt \(a\in D_p^S\).
Weil \(a\) beliebig war, gilt \(A\subseteq D_p^S\). Da \(A\) schon
nichtleer war, ist \(A\in\PowNE{(D_p^S)}\).

\medskip\noindent\textbf{Von den Teilmengen von \(D_p^S\) zu \(L_p^S\).}
Sei umgekehrt \(\varnothing\ne A\subseteq D_p^S\).
Für jedes \(a\in A\) gilt dann \(\{a\}\preceq_S B\).
Wir wollen daraus \(A\preceq_S B\) gewinnen. Gilt \(BX=\{0\}\),
so liefert der Vergleich jeder Einermenge \(\{a\}\) mit \(B\)
die Gleichheit \(\{a\}X=\{0\}\). Also ist \(ax=0\) für
\emph{alle} \(a\in A\) und \(x\in X\). Damit enthält \(AX\) nur
den Wert \(0\). Da \(A\) und \(X\) nichtleer sind, gibt es wenigstens
ein solches Produkt; somit ist \(AX=\{0\}\).

Gilt stattdessen \(XB=\{0\}\), so ist für jedes \(a\in A\)
auch \(X\{a\}=\{0\}\). Wieder sind alle Produkte \(xa\) gleich \(0\),
und wieder sorgt die Nichtleerheit für \(XA=\{0\}\).
Dies beweist \(A\preceq_S B\), also \(A\in L_p^S\).

Der Beweis verwendet nur die Nulltests und ein Nullelement.
Da \(T\) ebenfalls ein Nullelement besitzt, gilt derselbe Beweis
mit \(T\) anstelle von \(S\). Für jedes \(r\in Q_T\) erhalten wir
also auch \(L_r^T=\PowNE{(D_r^T)}\).

\begin{NPReadingExample}{Dreierbeispiel: Die Potenzmengenidentität}
Im Beispiel ist die Gleichheit unmittelbar sichtbar:
\(L_{p_0}^S\) besteht aus genau den drei nichtleeren Teilmengen von
\(\{0,c\}\), und \(L_{p_1}^S\) aus genau den sieben nichtleeren
Teilmengen von \(S\). Der allgemeine Beweis zeigt, dass dies kein
Zufall der kleinen Tafel ist. Er verwendet hier nur Nulltests
und gilt für jede Halbgruppe mit Null.
\end{NPReadingExample}

\section*{Wie \(\Phi\) die Nulltests auf die andere Seite überträgt}

In der Potenzhalbgruppe von \(S\) ist \(\{0\}\) ein Nullelement:
Mit jeder nichtleeren Menge ist ihr Produkt auf beiden Seiten wieder
\(\{0\}\). Entsprechend ist \(\{0_T\}\) eine Null der Potenzhalbgruppe
von \(T\). Ein Nullelement ist eindeutig, denn für zwei Nullen
\(z,z'\) wäre \(zz'=z\) und zugleich \(zz'=z'\).

Das Bild \(\Phi(\{0\})\) ist ebenfalls eine Null: Jede Zielmenge
hat die Form \(\Phi(X)\), und
\[
 \Phi(\{0\})\Phi(X)=\Phi(\{0\}X)=\Phi(\{0\}).
\]
Auf der anderen Seite der Multiplikation gilt dasselbe.
Wegen der Eindeutigkeit folgt \(\Phi(\{0\})=\{0_T\}\).
Zusammen mit der Produkterhaltung und der Injektivität von \(\Phi\)
erhalten wir daher
\[
 AX=\{0\}\quad\Longleftrightarrow\quad
 \Phi(A)\Phi(X)=\{0_T\},
\]
und ebenso für \(XA\). Die Rückrichtung folgt, weil gleiche Bilder
unter einer injektiven Abbildung gleiche Urbilder haben.
Beim Übertragen wechseln beide Mengen: Aus der zu prüfenden Menge
\(A\) wird \(\Phi(A)\), aus der Testmenge \(X\) wird \(\Phi(X)\).
Verglichen werden also entsprechende Tests in den beiden Halbgruppen.

Wir zeigen jetzt ausdrücklich
\[
 A\preceq_S B\quad\Longleftrightarrow\quad
 \Phi(A)\preceq_T\Phi(B).
\]
Sei zunächst \(A\preceq_S B\), und sei \(Y\) irgendeine nichtleere
Teilmenge von \(T\) mit \(\Phi(B)Y=\{0_T\}\). Weil \(\Phi\)
surjektiv ist, gibt es \(X\) mit \(Y=\Phi(X)\). Wir können den
Test nun hinüber und wieder zurück übertragen:
\[
 \begin{aligned}
 \Phi(B)Y=\{0_T\}
 &\ \Longrightarrow\ BX=\{0\}\\
 &\ \Longrightarrow\ AX=\{0\}\\
 &\ \Longrightarrow\ \Phi(A)Y=\{0_T\}.
 \end{aligned}
\]
Die mittlere Implikation ist genau die Voraussetzung \(A\preceq_S B\).
Aus \(Y\Phi(B)=\{0_T\}\) folgt auf dieselbe Weise
\(Y\Phi(A)=\{0_T\}\). Also gilt der Vergleich auf der Zielseite.
Für die umgekehrte Richtung wenden wir dieselbe Überlegung auf
den inversen Isomorphismus \(\Phi^{-1}\) an. Es gehen also weder
Tests noch Vergleichsinformationen verloren.

\section*{Profile und Mengenfamilien abbilden}

Aus dem Testvergleich gewinnen wir eine neue Abbildung
\[
 \psi\colon Q_S\longrightarrow Q_T,\qquad
 \psi([A]_S)=[\Phi(A)]_T.
\]
Man nimmt eine Menge \(A\) aus einem Profil, bildet sie mit \(\Phi\)
ab und betrachtet dann das Profil der Bildmenge. Das ist eine
Zuordnung von \emph{Profilen}; \(\psi\) ist weder \(\Phi\) noch die
noch gesuchte Abbildung \(f\) der Grundelemente.

Die Vorschrift hängt nicht von der Wahl von \(A\) ab. Haben \(A\)
und \(B\) dasselbe Profil, gelten beide Präordnungsvergleiche zwischen
ihnen. Nach dem vorigen Abschnitt gelten dann auch beide Vergleiche
zwischen \(\Phi(A)\) und \(\Phi(B)\); also haben die Bilder dasselbe
Profil. Umgekehrt erzwingt die Gleichheit der Bildprofile auch die
Gleichheit der Ausgangsprofile. Das zeigt die Injektivität von \(\psi\).
Jedes Zielprofil hat einen Vertreter \(Y=\Phi(A)\), also ist \(\psi\)
auch surjektiv. Schließlich gilt
\[
 [A]_S\le[B]_S
 \quad\Longleftrightarrow\quad
 \psi([A]_S)\le\psi([B]_S).
\]
\(\psi\) ist damit ein Ordnungsisomorphismus: eine Bijektion, welche
die Ordnung in beiden Richtungen erhält.

Für \(p\in Q_S\) gilt
\[
 \begin{aligned}
 A\in L_p^S
 &\quad\Longleftrightarrow\quad[A]_S\le p\\
 &\quad\Longleftrightarrow\quad[\Phi(A)]_T\le\psi(p)\\
 &\quad\Longleftrightarrow\quad\Phi(A)\in L_{\psi(p)}^T.
 \end{aligned}
\]
Damit bildet \(\Phi\) jede Menge aus \(L_p^S\) in die richtige
Zielfamilie ab. Verschiedene Mengen aus \(L_p^S\) behalten verschiedene
Bilder, weil \(\Phi\) injektiv ist. Jede Menge aus
\(L_{\psi(p)}^T\) hat wegen der Surjektivität von \(\Phi\) ein
Urbild, das nach der obigen Äquivalenz zu \(L_p^S\) gehört.
Damit ist die Einschränkung von \(\Phi\) eine Bijektion
\[
 \Phi|_{L_p^S}\colon L_p^S\longrightarrow L_{\psi(p)}^T.
\]
Die Abbildungsregel bleibt also unverändert; wir erlauben lediglich
nur noch Eingaben aus \(L_p^S\). Die Elemente dieser Familien
sind weiterhin nichtleere Teilmengen. Als Gleichheit von Bildfamilien lautet
die Aussage \(\Phi[L_p^S]=L_{\psi(p)}^T\). Sie ist keine Aussage
über eine Abbildung der Grundelemente in \(D_p^S\).

\section*{Erst zählen wir die Grundelemente bis zu einem Profil}

Aus der bewiesenen Bijektion folgt \(|L_p^S|=|L_{\psi(p)}^T|\).
Andererseits haben wir auf beiden Seiten gezeigt, dass \(L\) jeweils
die gesamte nichtleere Potenzmenge von \(D\) ist. Also besitzen
\(D_p^S\) und \(D_{\psi(p)}^T\) gleich viele nichtleere Teilmengen.
Genau hier wird die zuvor bewiesene Gleichheit
\(L_p^S=\PowNE{(D_p^S)}\) gebraucht: Sie sichert, dass wir
\emph{alle} nichtleeren Teilmengen zählen. Aus der Größe einer
beliebig ausgewählten Mengenfamilie könnten wir diesen Schluss
nicht ziehen.
Warum haben sie deshalb auch gleich viele Elemente?

Eine endliche Menge mit \(n\) Elementen hat genau \(2^n-1\)
nichtleere Teilmengen. Beim Zusammenstellen einer Teilmenge gibt es
für jedes Element zwei Möglichkeiten: aufnehmen oder nicht aufnehmen.
Von den \(2^n\) Möglichkeiten wird anschließend die leere Menge
weggelassen. Die Zahlen \(2^n-1\) wachsen streng mit \(n\), denn
\[
 2^{n+1}-1=(2^n-1)+2^n>2^n-1.
\]
Bei jedem zusätzlichen Grundelement steigt die Teilmengenanzahl
also tatsächlich an. Deshalb können zwei verschiedene endliche
Elementzahlen nicht dieselbe Anzahl nichtleerer Teilmengen ergeben.
Der einschlägige Satz aus
Band~20 ist \FormulaRefAuto{FiniteNonemptyPowerSetReflectsCardinality}[thm].
Er ergibt
\[
 \boxed{|D_p^S|=|D_{\psi(p)}^T|.}
\]
Wir verwenden hier keine Logarithmusfunktion, sondern die Eindeutigkeit
der endlichen Elementzahl bei vorgegebener Potenzmengengröße.

\begin{NPReadingExample}{Dreierbeispiel: Die Zielanzahlen Schritt für Schritt}
Für das untere Profil wissen wir jetzt
\[
 \bigl|\PowNE{(D_{\psi(p_0)}^T)}\bigr|
   =|L_{\psi(p_0)}^T|=|L_{p_0}^S|=3.
\]
Die erste Gleichheit folgt aus der Potenzmengenidentität auf \(T\),
die zweite aus der eingeschränkten Bijektion \(\Phi\), die letzte
aus der Liste \(\{0\},\{c\},\{0,c\}\).
Weil \(2^2-1=3\) und keine andere Elementzahl dieselbe
Teilmengenanzahl liefert, hat \(D_{\psi(p_0)}^T\) genau zwei Elemente.

Für das obere Profil gilt entsprechend
\[
 \bigl|\PowNE{(D_{\psi(p_1)}^T)}\bigr|
   =|L_{\psi(p_1)}^T|=|L_{p_1}^S|=7.
\]
Hier werden alle sieben nichtleeren Teilmengen von \(S\) gezählt.
Aus \(2^3-1=7\) folgt daher, dass \(D_{\psi(p_1)}^T\) genau
drei Grundelemente hat. Die \(3\) in der ersten Rechnung zählt
\emph{Teilmengen}; diese letzte \(3\) zählt \emph{Grundelemente}.
\end{NPReadingExample}
Noch haben wir damit \emph{alle Grundelemente bis zu einem Profil}
gezählt. Als Nächstes trennen wir davon die Elemente ab, die schon
zu kleineren Profilen gehören.

\section*{Grundelemente mit genau diesem Profil zählen}

Jedes \(s\in S\) hat genau ein Einermengenprofil \([\{s\}]_S\).
Die Mengen \(C_q^S\) sind daher paarweise disjunkt und enthalten
zusammen alle Grundelemente. Sie heißen auch die \emph{Fasern}
der Zuordnung \(s\mapsto[\{s\}]_S\). Dabei darf eine Faser leer sein:
Ein Profil kann durch größere Teilmengen vertreten werden, ohne
irgendeine Einermenge zu enthalten.

Wie zuvor bedeutet \(q<p\) ausdrücklich:
\(q\le p\) und \(q\ne p\). Hier werden \emph{Profile} in der
Nulltestordnung verglichen, keine Zahlen und auch nicht ihre
Mächtigkeiten. Ein strikt kleineres Profil besteht mindestens alle
Tests von \(p\), stimmt aber nicht mit \(p\) überein.

Wir sammeln zunächst alle Grundelemente, deren Profil strikt kleiner
als \(p\) ist, in einer eigenen Menge:
\[
 E_p^S=\{s\in S\mid[\{s\}]_S<p\}.
\]
Dies ist genau der Teil von \(D_p^S\), der nicht zu \(C_p^S\) gehört.
Denn für \(s\in D_p^S\) gilt \([\{s\}]_S\le p\); liegt \(s\)
nicht in \(C_p^S\), ist dieses Profil zusätzlich von \(p\) verschieden.
Umgekehrt erfüllt jedes \(s\) mit \([\{s\}]_S<p\) genau diese
beiden Bedingungen. Somit gilt
\[
 E_p^S=D_p^S\setminus C_p^S,\qquad
 D_p^S=C_p^S\mathbin{\dot\cup}E_p^S.
\]
Der Punkt über dem Vereinigungszeichen bedeutet: Die beiden Teile
überschneiden sich nicht. Sie teilen \(D_p^S\) vollständig in
\emph{genau dieses Profil} und \emph{strikt kleinere Profile} auf.

Sortieren wir die Elemente von \(E_p^S\) zusätzlich nach ihrem
jeweiligen Profil, ergibt sich eine zweite Beschreibung:
\[
 E_p^S=\bigcup_{\substack{q\in Q_S\\q<p}}C_q^S.
\]
Ist nämlich \(s\in E_p^S\), so wählen wir sein eindeutig bestimmtes
Profil \(q=[\{s\}]_S\). Es erfüllt \(q<p\), und \(s\in C_q^S\)
zeigt die Zugehörigkeit zur Vereinigung. Gehört umgekehrt \(s\)
zu dieser Vereinigung, liegt es in einer Faser \(C_q^S\) mit \(q<p\).
Sein Profil ist also \(q<p\), genau die Bedingung für \(s\in E_p^S\).
Der Vereinigungsindex durchläuft nur die strikt kleineren Profile;
unvergleichbare Profile werden nicht mitgezählt.

Auf der Zielseite gilt dieselbe Beschreibung. Für ein beliebiges
\(r\in Q_T\) setzen wir
\[
 \begin{aligned}
 E_r^T&=\{t\in T\mid[\{t\}]_T<r\}\\
      &=D_r^T\setminus C_r^T
       =\bigcup_{\substack{q\in Q_T\\q<r}}C_q^T.
 \end{aligned}
\]
Die Begründung ist dieselbe elementweise Unterscheidung: Ein Profil
höchstens \(r\) ist entweder gleich \(r\) oder strikt kleiner.
Der Buchstabe \(r\) bezeichnet hier ein beliebiges Zielprofil.
Beim Vergleich mit \(S\) werden wir speziell \(r=\psi(p)\) einsetzen.

\begin{NPReadingExample}{Dreierbeispiel: Was genau wird abgezogen?}
Strikt unterhalb von \(p_0\) gibt es kein Profil. Deshalb ist die Vereinigung
über \(q<p_0\) leer: \(E_{p_0}^S=\varnothing\).
Die beiden Elemente von \(D_{p_0}^S=\{0,c\}\) gehören somit beide
zu \(C_{p_0}^S\).
Strikt unterhalb von \(p_1\) liegt genau \(p_0\). Daher gilt
\[
 \underbrace{D_{p_1}^S}_{\{0,c,a\}}
 =\underbrace{C_{p_1}^S}_{\{a\}}
  \mathbin{\dot\cup}\underbrace{E_{p_1}^S}_{C_{p_0}^S=\{0,c\}}.
\]
Auf der Zielseite gibt es ebenfalls genau die beiden Profile
\(\psi(p_0)<\psi(p_1)\). Bereits bewiesen sind die Größen zwei und
drei der zugehörigen \(D\)-Mengen. Strikt unter \(\psi(p_0)\) liegt kein
Profil, also hat seine \(C\)-Faser zwei Elemente. Diese beiden
Elemente bilden genau \(E_{\psi(p_1)}^T\). Von den drei Elementen
in \(D_{\psi(p_1)}^T\) bleibt somit eines für
\(C_{\psi(p_1)}^T\). Auf beiden Seiten werden also von drei
Grundelementen zwei entfernt; genau ein Grundelement bleibt.
Die Rechnung \(3-2=1\) betrifft die \(D\)-, \(E\)- und \(C\)-Mengen,
nicht die Anzahlen von Teilmengen in den \(L\)-Familien.
\end{NPReadingExample}

Im allgemeinen Fall können Profile unvergleichbar sein. Wir dürfen
sie also nicht stillschweigend als eine Kette behandeln. Stattdessen
verwenden wir die Minimalität in einer endlichen Halbordnung.
Wir beweisen damit
\[
 \boxed{|C_p^S|=|C_{\psi(p)}^T|\quad\text{für jedes }p\in Q_S.}
\]
Angenommen, es gäbe Profile, bei denen die beiden Fasergrößen
verschieden sind. Da \(S\) endlich ist, ist auch seine nichtleere
Potenzmenge endlich; somit hat sie nur endlich viele Profile.
Wir beginnen bei irgendeinem Profil, bei dem sich die beiden
Fasergrößen unterscheiden. Gibt es ein strikt kleineres Profil,
bei dem sie ebenfalls verschieden sind, gehen wir zu einem solchen über. Diesen Schritt
wiederholen wir, solange es möglich ist. Dabei kehren wir zu keinem
früheren Profil zurück, denn jeder Schritt führt strikt nach unten.
Da es nur endlich viele Profile gibt, muss der Abstieg enden.
Das Endprofil nennen wir \(p\). Es ist ein \emph{minimales Gegenbeispiel}:
Bei \(p\) sollen die beiden Fasergrößen verschieden sein, bei jedem
strikt kleineren Profil stimmen sie dagegen überein.
Wir benötigen kein kleinstes Profil aller
Gegenbeispiele; mit \(p\) unvergleichbare Profile sind hier unerheblich.
Für jedes \(q<p\) wissen wir also
\(\lvert C_q^S\rvert=\lvert C_{\psi(q)}^T\rvert\).

Zwischen diesen entsprechenden kleineren Fasern können wir Bijektionen
wählen. Zusammen ergeben sie eine Bijektion
\(g_p\colon E_p^S\to E_{\psi(p)}^T\).
Diese Hilfsabbildung vergleicht vorerst nur die Elemente mit Profil
strikt unterhalb des fest gewählten \(p\), um ihre Anzahlen abzuziehen.
Gibt es kein Profil strikt unter \(p\), gibt es wegen des
Ordnungsisomorphismus auch keines strikt unter \(\psi(p)\).
Beide \(E\)-Mengen sind dann leer, und \(g_p\) ist die leere Bijektion.
\par\Needspace{6\baselineskip}
Für ihre Bijektivität sind drei Punkte entscheidend.
Erstens gehört jedes Ausgangselement zu genau einer Faser, sodass ihm
nur ein Wert zugeordnet wird. Zweitens führen verschiedene Ausgangsprofile
zu verschiedenen Zielprofilen, weil \(\psi\) injektiv ist; die Bildfasern
überschneiden sich daher nicht.

\par\Needspace{6\baselineskip}
Drittens entsteht jede Faser
strikt unterhalb von \(\psi(p)\) auf diese Weise: Aus dem Ordnungsisomorphismus
folgt
\[
 q<p\quad\Longleftrightarrow\quad\psi(q)<\psi(p),
\]
und jedes Zielprofil hat ein Urbild. Es fehlt also kein Element
von \(E_{\psi(p)}^T\). Die Wahl und das Zusammenfügen solcher
Bijektionen sind durch den Faserbijektionssatz aus Band~09 abgesichert.
Hier werden nur endlich viele Bijektionen benötigt.

Nun vergleichen wir die beiden Zerlegungen
\[
 \begin{aligned}
 D_p^S&=C_p^S\mathbin{\dot\cup}E_p^S,\\
 D_{\psi(p)}^T&=C_{\psi(p)}^T\mathbin{\dot\cup}E_{\psi(p)}^T.
 \end{aligned}
\]
Die beiden Gesamtmengen sind gleich groß, wie der vorherige Abschnitt
bewiesen hat. Auch die beiden \(E\)-Teile sind gleich groß, wie die
eben konstruierte Bijektion zeigt. In zwei gleich großen endlichen
Mengen bleiben nach dem Entfernen gleich großer Teile gleich große
Reste. Im Skript ist dies der Komplementsatz
\FormulaRefAuto{FiniteComplementBijection}[thm] aus Band~20.
Seine Voraussetzungen sind hier die Endlichkeit der \(D\)-Mengen,
eine Bijektion zwischen ihnen und eine Bijektion zwischen ihren
\(E\)-Teilen. Ihre Komplemente sind genau die beiden \(C\)-Fasern.
Folglich sind auch diese gleich groß.

Dabei behaupten wir nicht, dass eine beliebig gewählte Bijektion
der Gesamtmengen schon die \(E\)-Teile passend zuordnet.
Gerade der Komplementsatz erlaubt es, aus den beiden vorhandenen
Bijektionen eine Bijektion der verbleibenden Teile zu gewinnen.
Das Ergebnis widerspricht der Wahl von \(p\) als Gegenbeispiel und
beweist die Gleichheit für alle Profile, auch für solche mit leerer Faser.

Der Satz \FormulaRefAuto{FinitePosetCumulativeTransportBijection}[thm]
bündelt diese Argumentation in Band~20. Sein Inhalt ist hier konkret:
Wenn alle Anzahlen \emph{bis zu} einem Profil übereinstimmen, müssen
auch die Anzahlen \emph{genau bei} jedem Profil übereinstimmen.

\section*{Warum wir die Bilder von \(0\) und \(c\) vorschreiben dürfen}

Bevor wir die Faserbijektionen zur gesuchten Abbildung \(f\)
zusammensetzen, sichern wir die beiden nötigen Vorgaben ab:
\[
 f(0)=0_T,\qquad f(c)=c_T.
\]

Die bewiesene Gleichheit der Fasergrößen erlaubt innerhalb jeder
Faser die Wahl einer Bijektion. Wir prüfen zuerst, ob die
gewünschten Bilder überhaupt in den richtigen Fasern liegen.

Für \(0\) folgt das aus \(\Phi(\{0\})=\{0_T\}\):
\[
 \psi([\{0\}]_S)=[\{0_T\}]_T.
\]
Für \(c\) benutzen wir die bewiesene Bijektion
\(\Phi_2\colon\mathcal M_S\to\mathcal M_T\) zwischen den beiden
Familien aus je drei Produktmengen. Nachdem \(\{0\}\) auf \(\{0_T\}\)
abgebildet ist, bleiben als Bilder von \(\{c\}\) nur
\(\{c_T\}\) oder \(\{0_T,c_T\}\) übrig.

Diese beiden Mengen haben dasselbe Nullproduktprofil.
Für jeden rechten Test gilt nämlich
\[
 \{0_T,c_T\}X=\{0_T\}\cup\bigl(\{c_T\}X\bigr).
\]
Die rechte Seite ist genau dann \(\{0_T\}\), wenn schon
\(\{c_T\}X=\{0_T\}\) ist. Hier verwenden wir wieder, dass
\(\{c_T\}X\) nichtleer ist. Für linke Tests gilt dieselbe Überlegung.
Das Hinzufügen einer Null ändert also keinen Nulltest. Folglich ist
\[
 \psi([\{c\}]_S)=[\{c_T\}]_T.
\]
Beide vorgeschriebenen Bilder gehören damit zu den passenden Zielfasern.

Nun müssen wir nur noch die Bijektionen in diesen Fasern passend wählen.
Warum ist das möglich? Sei \(g\colon C\to C'\) eine vorhandene Bijektion,
und soll ein bestimmtes \(x\in C\) auf \(y\in C'\) abgebildet werden.
Falls \(g(x)=y\), ist nichts zu tun. Andernfalls vertauschen wir in
der Zielmenge die beiden Werte \(g(x)\) und \(y\) und lassen alle
anderen Werte fest. Diese Vertauschung ist selbst eine Bijektion.
Nach ihr hat \(x\) das gewünschte Bild, und die gesamte Abbildung
ist weiterhin bijektiv.

Liegen \(0\) und \(c\) in verschiedenen Fasern, führen wir diesen
Schritt unabhängig in jeder der beiden Fasern aus.
Liegen sie in derselben Faser, stellen wir zuerst das Bild von \(0\)
auf \(0_T\) ein. Weil die korrigierte Abbildung injektiv ist und
\(c\ne0\), ist ihr Wert bei \(c\) nun verschieden von \(0_T\).
Auch \(c_T\) ist verschieden von \(0_T\).
Wenn wir anschließend den bisherigen Wert bei \(c\) mit \(c_T\)
vertauschen, ist \(0_T\) daher an dieser zweiten Vertauschung nicht
beteiligt. Das erste vorgeschriebene Bild bleibt richtig.
Dies ist der Zweipunktsatz
\FormulaRefAuto{TwoPointBijectionExists}[thm] aus Band~08.

\begin{NPReadingExample}{Dreierbeispiel: Die beiden vorgeschriebenen Bilder}
Im Dreierbeispiel liegen \(0\) und \(c\) tatsächlich in derselben
Faser \(\{0,c\}\). Beide ergeben mit jedem Element Nullprodukte;
Nulltests allein unterscheiden ihre Rollen nicht.
Gerade deshalb ist dieser Fall kein nebensächlicher Sonderfall.
Die passende Wahl schickt die beiden verschiedenen Elemente auf
\(0_T\) und \(c_T\).
\end{NPReadingExample}

Damit können wir die Faserbijektionen von Anfang an so wählen, dass
sie \(0\) auf \(0_T\) und \(c\) auf \(c_T\) abbilden.
Die übrigen Faserbijektionen dürfen beliebig gewählt werden.
Zur formalen Profilverträglichkeit der markierten Werte:
\NullproductsProofLink{NPMarkedProfileCompatibility}.

\section*{Jetzt ordnen wir die Grundelemente einander zu}

Wir können nun für jedes Profil \(p\in Q_S\) eine Bijektion
\[
 f_p\colon C_p^S\longrightarrow C_{\psi(p)}^T
\]
wählen. Wir treffen die Wahl mit den gerade abgesicherten Vorgaben
für \(0\) und \(c\). Sind beide Fasern leer, nehmen wir die leere
Abbildung. Nun setzen wir die Bijektionen über \emph{alle} Profile
zusammen; ihre Definitionsbereiche zerlegen ganz \(S\).
Für \(s\in S\) setzen wir \(f(s)=f_p(s)\), wobei
\(p=[\{s\}]_S\) das eindeutig bestimmte Profil von \(s\) ist.
Diese Vorschrift gibt jedem Grundelement genau ein Bild und erfüllt
bereits \(f(0)=0_T\) und \(f(c)=c_T\).

Warum ist \(f\) bijektiv? Innerhalb einer Faser ist \(f_p\) injektiv.
Elemente aus verschiedenen Fasern gelangen in verschiedene Bildfasern,
können also ebenfalls nicht dasselbe Bild haben. Das beweist die
Injektivität. Jedes \(t\in T\) liegt in einer Faser \(C_r^T\).
Weil \(\psi\) surjektiv ist, gilt \(r=\psi(p)\) für ein \(p\), und
weil \(f_p\) surjektiv ist, hat \(t\) ein Urbild in \(C_p^S\).
Das beweist die Surjektivität. Genau diese Art des Zusammenfügens ist
in \FormulaRefAuto{FiberwiseBijectionExists}[thm] aus Band~09 formuliert.

Die Konstruktion legt außerdem fest, welches Profil die Einermenge
eines Bildes besitzt:
\[
 [\{f(s)\}]_T=\psi([\{s\}]_S)=[\Phi(\{s\})]_T.
\]
Die erste Gleichheit folgt aus \(f(s)\in C_{\psi(p)}^T\), die zweite
ist die Definition von \(\psi\). Das besagt \emph{nicht}, dass
\(\{f(s)\}=\Phi(\{s\})\) ist. Die beiden Mengen haben dieselben
Nulltests; \(\Phi(\{s\})\) darf durchaus mehrere Elemente enthalten.

\section*{Warum diese Zuordnung die Nullprodukte erhält}

Zunächst klären wir einen Austauschschritt. Haben \(A,A'\) dasselbe
Profil und ebenso \(B,B'\), dann gilt
\[
 AB=\{0\}\quad\Longleftrightarrow\quad A'B'=\{0\}.
\]
Denn wir dürfen zuerst den rechten Test mit \(X=B\) für die
profilgleichen Mengen \(A,A'\) einsetzen. Das gibt
\(AB=\{0\}\Longleftrightarrow A'B=\{0\}\).
Danach setzen wir den linken Test mit \(X=A'\) für \(B,B'\) ein.
Das gibt \(A'B=\{0\}\Longleftrightarrow A'B'=\{0\}\).
So wird ein Faktor nach dem anderen ausgetauscht. Dasselbe Argument
gilt mit \(0_T\) auf der Zielseite.

Für beliebige \(s,t\in S\) haben \(\Phi(\{s\})\) und \(\{f(s)\}\)
dasselbe Profil. Ebenso gilt dies für \(\Phi(\{t\})\) und \(\{f(t)\}\).
Auf diese beiden Paare wenden wir den Austauschschritt an.
Wir erhalten die vollständige Kette
\[
 \begin{aligned}
 st=0
 &\quad\Longleftrightarrow\quad\{s\}\{t\}=\{0\}\\
 &\quad\Longleftrightarrow\quad
       \Phi(\{s\})\Phi(\{t\})=\{0_T\}\\
 &\quad\Longleftrightarrow\quad
       \{f(s)\}\Phi(\{t\})=\{0_T\}\\
 &\quad\Longleftrightarrow\quad
       \{f(s)\}\{f(t)\}=\{0_T\}\\
 &\quad\Longleftrightarrow\quad f(s)f(t)=0_T.
 \end{aligned}
\]
Am Anfang und am Ende wird ein Einermengenprodukt als Elementprodukt
gelesen. Der Übergang mit \(\Phi\) verwendet dessen Produkterhaltung,
Injektivität und \(\Phi(\{0\})=\{0_T\}\).
Die beiden mittleren Austauschschritte verwenden die eben erklärten
Nulltests. Damit kennen wir für jedes Paar von Grundelementen genau,
ob sein Produkt Null ist.

Die eben bewiesene Erhaltung der Nullprodukte benutzt nur die
Profilgleichheiten. Sie gilt für jede Bijektion, die entsprechende
Fasern zuordnet. Die zusätzlich festgelegten Bilder von \(0\) und
\(c\) ermöglichen jetzt den letzten Schritt: Aus dem Nullverhalten
bestimmen wir den genauen Wert jedes Produkts.

\par\Needspace{12\baselineskip}
\section*{Jetzt sind sämtliche Produkte bestimmt}

Seien \(x,y\in S\) beliebig. Es gibt genau zwei Fälle, denn
\(xy\in S^2=\{0,c\}\).

Ist \(xy=0\), so liefert die Nullprodukterhaltung
\(f(x)f(y)=0_T\). Wegen der vorgeschriebenen Zuordnung von \(0\)
ist gleichzeitig \(f(xy)=f(0)=0_T\). Also gilt
\[
 f(xy)=f(x)f(y).
\]
Ist \(xy=c\ne0\), so muss \(f(x)f(y)\ne0_T\) sein: Andernfalls
würde die \emph{Rückrichtung} der Nullprodukterhaltung \(xy=0\)
erzwingen. Aber jedes Produkt in \(T\) liegt in \(\{0_T,c_T\}\).
Als einziger Nichtnullwert bleibt deshalb \(c_T\).
Wegen \(f(c)=c_T\) folgt auch hier
\[
 f(x)f(y)=c_T=f(c)=f(xy).
\]
Damit erhält \(f\) für jedes Paar die Multiplikation.
Da \(f\) bereits bijektiv ist, ist es der gesuchte Halbgruppenisomorphismus.
Das beweist den Hauptsatz. Zur formalen Schlussableitung:
\NullproductsProofLink{NPReconstructionConclusion}.

\begin{NPReadingExample}{Dreierbeispiel: Die rekonstruierte Multiplikation}
Im Dreierbeispiel lässt sich der Schluss noch einmal an der Tafel lesen.
Die Grundelemente \(0,c\) bestehen jeden Nulltest, das verbleibende
Element \(a\) nicht. Die Größen der entsprechenden Zielfasern sind
zwei und eins. Wir schicken \(0\) auf \(0_T\), \(c\) auf \(c_T\)
und \(a\) auf das einzige Element \(a_T\) der anderen Faser.
Alle Produkte mit \(0\) oder \(c\) bleiben Nullprodukte.
Weil \(a^2\) nicht Null ist, ist auch \(a_T^2\) nicht Null und
daher gleich \(c_T\). Damit ist die ursprüngliche Dreiertafel
vollständig zurückgewonnen.
\end{NPReadingExample}

\section*{Was die Nullproduktinformation bei drei Werten offenlässt}

Beim letzten Schritt war entscheidend, dass es genau einen
Nichtnullproduktwert gibt. Wenn es mehrere gibt, beantwortet ein
Nulltest nur die Frage, ob ein Produkt Null ist.
Er unterscheidet die verschiedenen Nichtnullwerte nicht.
Das folgende Beispiel macht diese Grenze konkret.

Wir definieren zwei Multiplikationen auf \(\{0,a,b,u,v\}\).
Bei beiden ist ein Produkt gleich \(0\), sobald wenigstens einer
seiner Faktoren zu \(\{0,u,v\}\) gehört.
Nur die vier Paare aus \(\{a,b\}\) liefern andere Werte.
Für sie setzen wir
\[
 \begin{array}{c|cc}
 \cdot&a&b\\\hline
 a&u&v\\
 b&v&u
 \end{array}
 \qquad\qquad
 \begin{array}{c|cc}
 \ast&a&b\\\hline
 a&u&u\\
 b&v&u
 \end{array}
\]
Beispielsweise ist links \(a\cdot b=v\), rechts aber \(a\ast b=u\).
Keines der vier aufgeführten Produkte ist Null.
Für alle nicht aufgeführten Paare sind beide Produkte Null.
Die beiden Operationen haben somit genau dieselbe Nullproduktrelation.
Auch ihre Produktmengen sind beide \(\{0,u,v\}\), denn alle drei
Werte kommen vor.

Beide Operationen sind assoziativ. Ein zuerst gebildetes Produkt
liegt immer in \(\{0,u,v\}\), und jedes Element dieser Menge ergibt
mit jedem weiteren Element das Produkt \(0\). Bei jedem dreifachen
Produkt führen deshalb beide Klammerungen zum Wert \(0\).
Trotzdem sind die Halbgruppen nicht isomorph.
Die linke Tafel ist kommutativ, einschließlich aller nicht
aufgeführten Nullprodukte. Rechts gilt dagegen
\(a\ast b=u\ne v=b\ast a\).
Ein Isomorphismus einer kommutativen Halbgruppe würde wegen der
Surjektivität alle Zielpaare kommutieren lassen; eine nichtkommutative
Halbgruppe kann also nicht dazu isomorph sein.

Dieses Beispiel zeigt, warum unsere konstruierte Nullproduktbijektion
allein bei drei verschiedenen Produktwerten nicht genügt.
Diese beiden Halbgruppen haben jeweils fünf Grundelemente;
die Zahl drei bezeichnet hier die Produktwerte \(0,u,v\).
Das Beispiel behauptet nicht, dass die
\emph{ganzen Potenzhalbgruppen} dieses Beispiels isomorph wären.
Tatsächlich sind sie es nicht: Links ist auch die Mengenmultiplikation
kommutativ. Rechts unterscheiden bereits die Einermengenprodukte
\(\{a\}\{b\}\) und \(\{b\}\{a\}\) die beiden Reihenfolgen.
Die gesamte Potenzhalbgruppe enthält also mehr Information als die
Nulltests, auf die wir uns im mittleren Teil beschränkt haben.

\section*{Wie die Teile des Beweises zusammenwirken}

Die Ausgangsfrage lautete, wie man von einer Multiplikation der
Teilmengen zur Multiplikation der Grundelemente zurückkommt.
Der Beweis braucht dafür keinen Nachweis, dass der gegebene
Isomorphismus \(\Phi\) Einermengen erhält. Er gewinnt stattdessen
genügend Information, um eine geeignete neue Abbildung \(f\)
zu konstruieren.

Zuerst sichert die Bijektion der Potenzmengen die Endlichkeit von \(T\).
Die drei möglichen Produktmengen liefern sein Zweierquadrat;
dessen Multiplikation liefert die Null.
Erst danach vergleichen wir dieselben
Nulltests in beiden Halbgruppen. Die Profile halten fest,
welche Teilmengen bei diesen Tests ununterscheidbar sind.
Ihre Ordnung hält zusätzlich fest, welche Mengen mindestens
alle Tests anderer Mengen bestehen.

Die Profile allein sagen noch nicht, wie viele Grundelemente es gibt.
Dafür werden die Familien \(L_p^S\) gebraucht.
Ihre Mitglieder sind genau alle nichtleeren Teilmengen von \(D_p^S\).
Aus ihrer Anzahl gewinnen wir daher die Anzahl der Grundelemente
bis zum Profil \(p\). Das Abtrennen der kleineren Fasern liefert
anschließend die Anzahl der Grundelemente mit genau diesem Profil.
Im Anfangsbeispiel ist das die leicht nachprüfbare Rechnung:
drei Teilmengen bedeuten zwei Grundelemente im unteren Profil;
sieben Teilmengen bedeuten insgesamt drei Grundelemente,
also bleibt ein weiteres für das obere Profil.
Die Minimalitätsargumentation rechtfertigt dieselbe Idee für jede
endliche Profilordnung, auch mit unvergleichbaren Profilen und
leeren Fasern.

\par\Needspace{4\baselineskip}
Erst jetzt wählen wir die Bijektionen zwischen den Grundelementfasern.
Die freie Wahl innerhalb der Fasern erlaubt, \(0\) und \(c\) bereits
dabei auf die richtigen Produktwerte zu schicken.
Die zusammengesetzte Bijektion erhält die Nullprodukte, weil deren
Wahrheitswert nur von den Profilen der beiden Faktoren abhängt.
Danach ist jede Produktgleichheit erzwungen:
Null bleibt Null, und für ein Nichtnullprodukt gibt es nur noch
eine Möglichkeit.

Hier wird auch sichtbar, was die Voraussetzungen leisten.
Die Endlichkeit trägt die Rückrechnung aus den Potenzmengengrößen
und die Minimalitätsargumentation.
Die zwei Produktwerte ermöglichen die Bestimmung der Zielstruktur
und schließlich die eindeutige Zuordnung jedes Nichtnullprodukts.
Der bewiesene Satz löst damit genau den Spezialfall
\(S^2=\{0,c\}\) mit \(c\ne0\).
Die allgemeine endliche Potenzhalbgruppenvermutung wird dadurch
nicht entschieden.

\section*{Einordnung, Geschichte und Quellen}

Der Hauptsatz steht in Band~37. Die allgemeinen Grundlagen zu
Präordnungen, endlichen Anzahlen und Faserbijektionen sind in den
jeweils genannten früheren Bänden bewiesen; Band~33 behandelt
die Nullproduktprofile. Die \NullproductsProofsLink{} führen die
zusätzlichen Ableitungen für den vorliegenden Rekonstruktionssatz
in der formalen Sprache des Skripts aus.
Die Links bei den einzelnen Beweisschritten führen zu den
entsprechenden Tabellen.

Die historische Frage lässt sich für das allgemeine Thema beantworten:
Das Isomorphieproblem für Potenzhalbgruppen wurde seit den 1960er Jahren
insbesondere von Tamura und Shafer untersucht; ihre Arbeit
\emph{Power semigroups} erschien 1967 in \emph{Mathematica Japonica}
12, S.~25--32. Für beliebige Halbgruppen ist die Rekonstruktion nicht
möglich, wie das in Band~28 behandelte Gegenbeispiel von Mogiljanskaja
zeigt. Die veröffentlichte Originalarbeit lautet
\href{https://doi.org/10.1007/BF02389140}
{\emph{Non-isomorphic semigroups with isomorphic semigroups of subsets}},
\emph{Semigroup Forum} 6 (1973), S.~330--333.

Bei der endlichen Frage muss man die Quellen sorgfältig unterscheiden.
Tamura kündigte 1987 in
\href{https://doi.org/10.1007/978-94-009-3839-7_22}
{\emph{On the Recent Results in the Study of Power Semigroups}}
eine positive Antwort für alle endlichen Halbgruppen ausdrücklich
ohne Beweis an. Liu und Tringali bezeichnen dieselbe allgemeine
endliche Frage noch in der Einleitung ihrer Arbeit
\href{https://arxiv.org/html/2606.01917v1}
{\emph{Power Semigroups and Two Rigidity Theorems for Groups}}
(1.~Juni 2026, Frage~1.1) als offen. Dort findet sich auch ein
Überblick über die historische Entwicklung.

Für den hier behandelten besonderen Satz mit \(S^2=\{0,c\}\)
konnte bei der Quellenprüfung vom 26.~September 2026 keine eindeutige
frühere Zuschreibung festgestellt werden. Das bedeutet weder, dass
der Satz neu ist, noch, dass er in der Literatur nicht vorkommt.
Die genannten Arbeiten belegen den historischen Zusammenhang; sie
werden hier nicht als Beleg für genau diesen Spezialfall angeführt.

\par\Needspace{4\baselineskip}
Ausgangspunkt der Ausarbeitung ist die vom Verfasser bereitgestellte
\href{https://chatgpt.com/share/6ab6cc27-72d8-83ed-b081-c0255d3a43fe}
{geteilte Unterhaltung über die Rekonstruktion mittels Nullproduktprofilen}.
Der Beweis wurde in die Definitionen und zuvor bewiesenen Aussagen
des Skripts eingeordnet. Mit dieser Herkunftsangabe ist keine Aussage
über die Priorität oder Neuheit des Spezialfalls verbunden.

\noindent\NullproductsMainLink{}\enspace\textperiodcentered\enspace\NullproductsProofsLink{}
